\begin{trapbox}
\textbf{\textcolor{CTrap}{概念陷阱：把圆的垂径公式生搬到椭圆}}

\begin{center}
\begin{tikzpicture}[font=\small]
 \path[use as bounding box] (-3,-2.9) rectangle (10,3.2);
 \node at (0,2.9) {圆：$OM\perp AB$};
 \draw[blue!60!black,thick] (0,0) circle (2);
 \draw[magenta!70!black,very thick] (-1.6,1.2)--(1.6,1.2);
 \draw[teal!70!black,thick] (-1.6,1.2)--(0,0)--(1.6,1.2);
 \draw[dashed,thick] (0,0)--(0,1.2);
 \draw (-.18,1.2)--(-.18,1.02)--(0,1.02);
 \node[below] at (0,0) {$O$}; \node[above] at (0,1.2) {$M$};
 \node[left] at (-1.6,1.2) {$A$}; \node[right] at (1.6,1.2) {$B$};
 \node at (0,-2.6) {$OA=OB=2$};
 \begin{scope}[shift={(6.8,0)}]
 \node at (0,2.9) {椭圆：不能照搬垂直关系};
 \draw[blue!60!black,thick] (0,0) ellipse (2.5 and 1.5);
 \coordinate (A) at (2.5,0); \coordinate (B) at (0,1.5);
 \coordinate (M) at (1.25,.75);
 \draw[magenta!70!black,very thick] (A)--(B);
 \draw[teal!70!black,thick] (A)--(0,0)--(B);
 \draw[dashed,thick] (0,0)--(M);
 \node[below] at (0,0) {$O$};
 \fill (M) circle (1.4pt);
 \draw[gray,thin] (M)--(2.2,1.7);
 \node[above] at (2.2,1.7) {$M$};
 \node[right] at (A) {$A$}; \node[above] at (B) {$B$};
 \node at (0,-2.1) {$OA=\frac52\ne\frac32=OB$};
 \node at (0,-2.6) {$\overrightarrow{OM}\cdot\overrightarrow{AB}=-2\ne0$};
 \end{scope}
\end{tikzpicture}
\end{center}


错误：对任意椭圆的斜弦仍假定“中心到弦中点连线垂直于弦”，
或写出 $2\sqrt{r^2-d^2}$ 却没有一个对所有曲线上点恒定的圆半径 $r$。
正确：\textbf{圆}有 $OA=OB=r$，才可使用垂径定理配合等腰三角形；
一般圆锥曲线的斜弦要回到两点坐标差与勾股定理。
\textbf{对照检查：}圆的 $2\sqrt{r^2-d^2}$ 和通用的
$\sqrt{1+k^2}|x_1-x_2|$，在\emph{同一条圆弦}上求出的值必须一致。

\medskip
\textbf{\textcolor{CTrap}{易错点一：平方开方忘记绝对值}}

错误：$\sqrt{(x_1-x_2)^2}=x_1-x_2$。
正确：$\boxed{\sqrt{(x_1-x_2)^2}=|x_1-x_2|}$。
\textbf{最小反例：}若 $x_1=1,x_2=4$，差为 $-3$，但距离不能为负数。

\medskip
\textbf{\textcolor{CTrap}{易错点二：把纵坐标式用于 $k=0$}}

错误：对于水平直线使用 $\sqrt{1+1/k^2}|y_1-y_2|$。
正确：$k=0$ 时 $y_1=y_2$，直接 $|AB|=|x_1-x_2|$。
\textbf{辨析：}纵坐标公式有 $1/k^2$，必须要求 $k\ne0$，不能把 $0\cdot\infty$ 当作普通乘法。

\medskip
\textbf{\textcolor{CTrap}{易错点三：竖直线硬套横坐标公式}}

错误：直线 $x=3$ 的两交点横坐标相同，就误算 $|AB|=0$。
正确：竖直线的 $k$ 不存在，弦长直接 $|y_1-y_2|$。
\textbf{最小反例：}圆 $x^2+y^2=4$ 与 $x=0$ 的交点是 $(0,2),(0,-2)$，弦长为 $4$ 而非 $0$。

\medskip
\textbf{\textcolor{CTrap}{易错点四：误用韦达差平方公式}}

错误：把 $|x_1-x_2|$ 写成 $\sqrt{(x_1+x_2)^2+4x_1x_2}$。
正确：$\boxed{(x_1-x_2)^2=(x_1+x_2)^2-4x_1x_2}$。
\textbf{核对法：}直接展开两侧：$x_1^2-2x_1x_2+x_2^2$。

\medskip
\textbf{\textcolor{CTrap}{易错点五：判别式大于零之前忘看二次项}}

只有当 $\alpha\ne0$，消元方程确实是
$\alpha x^2+\beta x+\gamma=0$ 时，
$\Delta>0$ 才意味着两个不同实根。
\textbf{最小反例：}$x^2-y^2=1$ 与 $y=x+1$ 联立后为 $-2x-2=0$，
实际只有一个交点；不能把它当作一个标准二次方程。

\medskip
\textbf{\textcolor{CTrap}{易错点六：判别式速算式漏掉 $|\alpha|$}}

正确：由二次方程得
\[
 |x_1-x_2|=\frac{\sqrt\Delta}{|\alpha|},\qquad
 |AB|=\frac{\sqrt{1+k^2}\sqrt\Delta}{|\alpha|}.
\]
\textbf{最小反例：}$2x^2-2=0$ 的两根 $\pm1$，根差为 $2$；
$\sqrt\Delta=\sqrt{16}=4$，必须再除以 $|2|$。

\medskip
\textbf{\textcolor{CTrap}{易错点七：混淆“有交点”与“有一条非退化弦”}}

切线只有一个重合交点（$\Delta=0$），对应的两点间距离为 $0$，
不能把它当作题目要求的由\textbf{两个不同点}构成的弦。
同时排查增根、分母限制以及消元是否等价。

\medskip
\textbf{\textcolor{CTrap}{一分钟自检}}

\textbf{先交点数，再看方向；选好坐标，再算根差；开方取正，分母验零。}
\end{trapbox}
